\section{Proposisi Majemuk}

Proposisi majemuk dibentuk dengan menggabungkan proposisi atomik menggunakan satu atau lebih operator logika \cite{rosen2019,forallx}. Dalam logika proposisional, proposisi majemuk disebut pula formula molekuler (molecular formula), sebagai lawan dari proposisi atomik yang tidak dapat diuraikan lebih lanjut oleh operator logika \cite{wiki_proplogic}. Nilai kebenaran proposisi majemuk ditentukan sepenuhnya oleh nilai kebenaran komponen-komponennya dan operator yang menghubungkannya.

Untuk mengevaluasi proposisi majemuk yang mengandung beberapa operator, diperlukan aturan prioritas (precedence). Urutan prioritas operator dari tertinggi ke terendah umumnya: negasi ($\neg$), konjungsi ($\wedge$), disjungsi ($\vee$), implikasi ($\rightarrow$), bikondisional ($\leftrightarrow$) \cite{levin_discrete,rosen2019}. Dengan demikian, $\neg p \wedge q$ dibaca sebagai $(\neg p) \wedge q$, bukan $\neg(p \wedge q)$. Tanda kurung dapat digunakan untuk mengelompokkan dan mengubah urutan evaluasi.

\begin{contoh}
Untuk proposisi $(p \wedge q) \rightarrow \neg r$: ``Jika $p$ dan $q$ benar, maka $r$ salah.'' Langkah evaluasi: (1) tentukan nilai $p \wedge q$; (2) tentukan nilai $\neg r$; (3) evaluasi implikasi antara hasil (1) dan (2). Contoh lain: $p \vee q \rightarrow r$ dibaca $(p \vee q) \rightarrow r$ karena disjungsi mempunyai prioritas lebih tinggi daripada implikasi.
\end{contoh}

Prinsip komposisi (principle of composition) menjamin bahwa formula well-formed dibangun secara rekursif dari atom-atom melalui aplikasi operator \cite{wiki_proplogic}. Artinya, setiap proposisi majemuk dapat dianalisis menjadi pohon sintaks: simpul-simpul daun adalah proposisi atomik, simpul-simpul dalam adalah operator, dan nilai kebenaran merambat dari daun ke akar sesuai definisi setiap operator \cite{forallx}.

Beberapa proposisi majemuk memiliki sifat khusus. \textbf{Tautologi} adalah proposisi yang selalu benar untuk semua interpretasi (semua baris dalam tabel kebenaran bernilai B). \textbf{Kontradiksi} adalah proposisi yang selalu salah. \textbf{Kontingensi} adalah proposisi yang kadang benar dan kadang salah bergantung pada interpretasi \cite{rosen2019,copi2014}. Dua proposisi majemuk dikatakan \textbf{ekuivalen secara logika} jika mereka memiliki tabel kebenaran yang identik.

\begin{table}[htbp]
\centering
\small
\begin{tabular}{|l|>{\raggedright\arraybackslash}p{4.2cm}|>{\raggedright\arraybackslash}p{4.5cm}|}
\hline
\textbf{Jenis} & \textbf{Definisi} & \textbf{Contoh formula} \\
\hline
Tautologi & Selalu benar (semua baris B) & $p \vee \neg p$, \ $p \rightarrow p$ \\
\hline
Kontradiksi & Selalu salah (semua baris S) & $p \wedge \neg p$ \\
\hline
Kontingensi & Ada baris B dan ada baris S & $p \wedge q$, \ $p \rightarrow q$ \\
\hline
\end{tabular}
\caption{Tautologi, kontradiksi, dan kontingensi}
\label{tab:tautologi-kontradiksi}
\end{table}

\begin{contoh}
\textbf{Evaluasi proposisi majemuk langkah per langkah}

Untuk $(p \wedge q) \rightarrow \neg r$, prioritas: kurung dulu $(p \wedge q)$ dan $\neg r$, lalu $\rightarrow$. Tabel kebenaran (kolom antara membantu):

\begin{center}
\small
\begin{tabular}{|c|c|c|c|c|c|}
\hline
$p$ & $q$ & $r$ & $p \wedge q$ & $\neg r$ & $(p \wedge q) \rightarrow \neg r$ \\
\hline
B & B & B & B & S & S \\
B & B & S & B & B & B \\
B & S & B & S & S & B \\
B & S & S & S & B & B \\
S & B & B & S & S & B \\
S & B & S & S & B & B \\
S & S & B & S & S & B \\
S & S & S & S & B & B \\
\hline
\end{tabular}
\end{center}
Ini kontingensi: ada baris yang B dan ada yang S.
\end{contoh}

\begin{contoh}
\textbf{Evaluasi proposisi majemuk dengan Python} \cite{python_docs_boolean}

Program berikut mengevaluasi $(p \wedge q) \vee \neg r$ untuk semua kombinasi nilai $p$, $q$, $r$ dan memeriksa apakah suatu formula selalu benar (tautologi):

\begin{lstlisting}[style=pythonstyle, caption={Evaluasi proposisi majemuk dan cek tautologi sederhana}]
from itertools import product

def eval_formula(p, q, r):
    """(p and q) or (not r)"""
    return (p and q) or (not r)

print("p    q    r    (p and q) or (not r)")
print("-" * 40)
for p, q, r in product([True, False], repeat=3):
    result = eval_formula(p, q, r)
    row = "B" if p else "S", "B" if q else "S", "B" if r else "S"
    print(f"{row[0]:4} {row[1]:4} {row[2]:4}  {'B' if result else 'S'}")

# Cek tautologi: p or not p
def is_tautology_neg_or_pos():
    return all((p or (not p)) for p in [True, False])
print("\nTautologi (p or not p):", is_tautology_neg_or_pos())  # True
\end{lstlisting}
\end{contoh}

Penerapan proposisi majemuk dalam pemrograman dan spesifikasi perangkat lunak sangat luas. Kondisi percabangan, loop, dan asersi sering kali berupa proposisi majemuk. Pemahaman prioritas operator dan cara menulis ulang proposisi (misalnya, hukum De Morgan: $\neg(p \wedge q) \equiv \neg p \vee \neg q$) membantu menyederhanakan kode dan menghindari kesalahan logika \cite{levin_discrete,stanford_intrologic}.
